\par\medskip
\noindent
14.2. \underline{Relaciones del homomorfismo}
$i:\underline M_K\longrightarrow G_K$ \underline{con la monodromía.}

Sea $n$ un entero $>0$. De la sucesión exacta (14.1.7), mediante la
sucesión exacta de los
$\underline{\operatorname{Ext}}^i(\mathbb Z/n\mathbb Z,-)$ en el
sitio de los espacios analíticos rígidos sobre $K$, se obtiene una
sucesión exacta canónica (teniendo en cuenta que
$\underline M_K^{\mathrm{an}}$ no tiene torsión y que
$G_K^{\mathrm{an}}$ es $n$-divisible):

\begin{equation*}
 0\longrightarrow{}_nG_K^{\mathrm{an}}
 \longrightarrow{}_nA_K^{\mathrm{an}}
 \longrightarrow\underline M_{K,n}^{\mathrm{an}}
 \longrightarrow0\qquad ,
\end{equation*}

\noindent
donde
$\underline M_{K,n}\overset{\mathrm{dfn}}{=}
\underline M_K/n\underline M_K$. Puesto que todos los términos que
aparecen en esta sucesión exacta proceden de esquemas en grupos finitos
sobre $K$, puede reescribirse como procedente de una sucesión exacta

\begin{equation}\tag{14.2.1}
 0\longrightarrow{}_nG_K\longrightarrow{}_nA_K
 \longrightarrow\underline M_{K,n}\longrightarrow0\qquad .
\end{equation}

\noindent
Cuando varía $n$, estas sucesiones exactas se comportan de la manera
esperada y, al pasar al límite inverso sobre las potencias de un número
primo $\ell$, se obtiene la sucesión exacta análoga

\begin{equation}\tag{14.2.2}
 0\longrightarrow T_\ell(G_K)\longrightarrow T_\ell(A_K)
 \longrightarrow\underline M_K\otimes_{\mathbb Z}\mathbb Z_\ell
 \longrightarrow0\qquad .
\end{equation}

\noindent
Pero ya se dispone de una sucesión exacta semejante por medios puramente
algebraicos, gracias a la primera fórmula (8.1.2) y a (8.1.7), que dan

\begin{equation}\tag{14.2.3}
 T_\ell(G_K)\xrightarrow{\ \sim\ }T_\ell(A_K)^{\mathrm{ef}}
 \subset T_\ell(A_K),\qquad
 T_\ell(A_K)/T_\ell(A_K)^{\mathrm{ef}}
 \simeq\underline M_K\otimes_{\mathbb Z}\mathbb Z_\ell\qquad .
\end{equation}

\noindent
Sentado esto, afirmo que \underline{los homomorfismos} (14.2.3)
\underline{procedentes de la teoría algebraica son precisamente los
deducidos de la sucesión exacta} (14.2.2) \underline{procedente de la
teoría trascendente}. Para el primer homomorfismo, que corresponde al
primer homomorfismo de (14.2.2), nuestra afirmación se sigue
inmediatamente, tras reducirnos a la situación semiestable, de la
conmutatividad de (14.1.11), pues las flechas verticales inducen allí
isomorfismos sobre los $T_\ell$, mientras que aplicar $T_\ell$ a las
flechas horizontales da, respectivamente, la flecha considerada en
(14.2.3) y la de (14.2.2). Admitiremos la compatibilidad afirmada para
los dos isomorfismos obtenidos entre
$T_\ell(A_K)/T_\ell(A_K)^{\mathrm{ef}}$ y
$\underline M_K\otimes_{\mathbb Z}\mathbb Z_\ell$. Obsérvese además
que \underline{esta compatibilidad} (para un $\ell$ dado)
\underline{basta ya para caracterizar de manera única, por medios
trascendentes, el homomorfismo} (14.1.5) \underline{que aparece en}
(14.1.7), interpretado como un isomorfismo de
$\underline M_K^{\mathrm{an}}$ sobre $\operatorname{Ker}e$.

También puede decirse que el homomorfismo $e$ de (14.1.7) identifica
los $T_\ell(A_K)/T_\ell(A_K)^{\mathrm{ef}}$ con los completados
$\ell$-ádicos $\underline R_K\otimes\mathbb Z_\ell$ de una única red
$\underline R_K=\operatorname{Ker}e$ sobre $K$, y que esta «estructura
entera» sobre el sistema precedente de grupos pro-$\ell$ es compatible,
en el sentido evidente, con la estructura entera procedente de las
fórmulas (8.1.7) (teorema de ortogonalidad 5.2). Esto define entonces
canónicamente un isomorfismo

\begin{equation*}
 i:\underline M_K\xrightarrow{\ \sim\ }\operatorname{Ker}e
 \qquad .
\end{equation*}

Puesto que el homomorfismo $i$ (14.1.5) permite reconstruir la
estructura de extensión usual sobre $T_\ell(A_K)$, a saber, (14.2.2),
y, por tanto, también su estructura de \underline{extensión mixta},
descrita en 9.6, mediante la sucesión exacta deducida de (14.1.1),

\begin{equation*}
\begin{array}{ccccccccc}
0&\longrightarrow&T_\ell(T_K)&\longrightarrow&T_\ell(G_K)
 &\longrightarrow&T_\ell(B_K)&\longrightarrow&0\\[5ex]
 &&\underline M'_K\otimes\mathbb Z_\ell(1)&&&&&&
\end{array}
\end{equation*}

\noindent
permite también, en el caso semiestable, reconstruir los
\underline{emparejamientos de monodromía} (9.1.2)

\begin{equation}\tag{14.2.4}
 u_\ell:\underline M\otimes\underline M'
 \longrightarrow\mathbb Z_\ell\qquad ,
\end{equation}

\noindent
que fueron definidos precisamente mediante extensiones mixtas en 9.6.
Limitémonos a indicar el resultado de traducir la construcción de 9.6.
Para ello, obsérvese que, en el caso semiestable, $p_K$ y $p'_K$
proceden de homomorfismos

\begin{equation*}
 p:\underline M\longrightarrow B,\qquad
 \sisi{\underline M'}{p':\underline M'}\longrightarrow B'\qquad ,
\end{equation*}

\noindent
donde $B,B'$ son esquemas abelianos sobre $S$, duales entre sí. Al
traer por $(p,p')$ se obtiene, por tanto, una biextensión $E$ de
$(\underline M,\sisi{M'}{\underline M'})$ por $\mathbb G_{mS}$ cuya
fibra genérica es la biextensión $E_K$ ya considerada en
\sisi{4.1}{14.1}.

Dicho esto, supongamos para simplificar que $\underline M$ y
$\underline M'$ son constantes, de valores $M$ y $M'$. Para cada par
$(m,m')\in M\times M'$, $E_{m,m'}$ es un torsor bajo
$\mathbb G_{mS}$ y corresponde, por tanto, a un módulo invertible
$L_{m,m'}$ sobre el anillo $V$ de $S$. Admite así una base $u_0$,
determinada salvo multiplicación por una unidad de $V$. En consecuencia,
se obtiene una aplicación canónica de «valoración»

\begin{equation*}
 \underline v:E_{m,m'}(K)\longrightarrow\mathbb Z
\end{equation*}

\noindent
que asocia a cada base $u$ de $L_{m,m'}\otimes_V K$, escrita
$u=\lambda u_0$ ($\lambda\in K^*$), el elemento

\begin{equation*}
 \underline v(u)=\underline v(\lambda)\in\mathbb Z\qquad ,
\end{equation*}

\noindent
donde $\underline v$ en el segundo miembro denota la valoración
asociada a $V$ (normalizada de la manera usual, de modo que su grupo de
valores sea $\mathbb Z$); el entero resultante es evidentemente
independiente de la elección de $u_0$. Si $\varphi$ es ahora una
trivialización de $E$, puede considerarse la aplicación correspondiente

\begin{equation}\tag{14.2.5}
 (m,m')\longmapsto\underline v\bigl(\varphi(m,m')\bigr):
 M\times M'\longrightarrow\mathbb Z\qquad .
\end{equation}

\noindent
Sentado esto, se encuentra que el \underline{homomorfismo de
monodromía} (14.2.4) \underline{toma sus valores en} $\mathbb Z$
\underline{y es precisamente el homomorfismo} (14.2.5), \underline{en
el que} $\varphi$ \underline{es la escisión de} $E_K$
\underline{asociada al homomorfismo} $i$ (14.1.5).

\medskip
\noindent
\underline{Observaciones 14.2.6.} En particular, se recupera (mediante
la teoría trascendente analítica rígida) el hecho de que el
emparejamiento de monodromía es entero e independiente de $\ell$
(10.4 a)), demostrado por otro método (también trascendente) en 12.5.
Según M.~Raynaud, la teoría analítica rígida permite recuperar asimismo
la afirmación de positividad 10.4 b) (que también está contenida en
12.5).

\medskip
\noindent
\underline{Observaciones 14.2.7.} Considérese
$\underline M_K\otimes_{\mathbb Z}\mathbb Q$ como una extensión de
$\underline M_K\otimes_{\mathbb Z}\mathbb Q/\mathbb Z$ por
$\underline M_K$. Mediante el homomorfismo (14.1.5)
$i:\underline M_K\longrightarrow G_K$, se deduce una extensión
$A_K^{\flat}$

\begin{equation}\tag{14.2.8}
 0\longrightarrow G_K\longrightarrow A_K^{\flat}
 \longrightarrow\underline M_K\otimes\mathbb Q/\mathbb Z
 \longrightarrow0\qquad .
\end{equation}

\noindent
Obsérvese además que, como $G_K$ es un grupo divisible, el dato de una
extensión semejante equivale al de una familia de extensiones de la
forma (14.2.2), a saber, las deducidas de (14.2.8) aplicando los
$T_\ell$. Puesto que $i$ es un monomorfismo, las extensiones en cuestión
que definen (14.2.8) son precisamente las extensiones (14.2.2)
procedentes de la sucesión exacta (14.2.1). Como estas últimas
extensiones están definidas por medios puramente algebraicos, lo mismo
sucede con la extensión (14.2.8). De hecho, la caracterización que
acabamos de dar muestra inmediatamente que esta extensión es
canónicamente isomorfa al grupo, denotado asimismo por $A_K^{\flat}$,
introducido en (8.2.3), en virtud del isomorfismo 11.11 del grupo de
componentes conexas de este último con
$\underline M_K\otimes\mathbb Q/\mathbb Z$.
